%----------------------------------------------------------------------
\section{Lagrangian Dual over the Resource Vector}
\label{sec:lagrangian}\footnote{Section-level Toulmin. \emph{Move}: formulate the planner's constrained-optimization problem as a linear program, derive its dual, and show how the dual prices give a direct diagnostic of binding resources. \emph{[DAF] comparison}: [DAF]'s $\kappa^\perp = \{0\}$ characterization of a saturated regime is the convex-optimization analogue of our dual-feasibility condition --- a regime is saturated exactly when no further resource can be spared. \emph{Residue instance}: the slack variable in the LP ($c - B$) is the budget's residue, and the dual variable (the $\lambda$) is the planner's shadow price for that residue.}
%----------------------------------------------------------------------

\subsection{The primal}

The planner's optimization problem is: find the cheapest
admissible plan $P$ for a workload $W$.

\begin{definition}[Primal plan selection]
  \label{def:primal}
  Given workload $W$, grammar $G$, budget vector
  $B \in \mathbb{N}^4$, and a scalar objective function
  $\phi : \mathbb{N}^4 \to \mathbb{R}$, the primal is
  \begin{equation}
    \min_{P \in \mathrm{plans}_G(W)} \phi(\bar{c}(P))
    \quad \text{s.t.} \quad \bar{c}(P) \leq B
    \text{ componentwise}.
  \end{equation}
\end{definition}

Choosing $\phi$ is a policy decision. Two natural
choices:

\begin{itemize}[nosep]
  \item $\phi(\bar{c}) = \bar{c}_{\LAT}$: minimize
    latency subject to memory budgets. Most common
    for latency-sensitive workloads.
  \item $\phi(\bar{c}) = \bar{c}_{\FLOPS}$ if we
    extend the resource vector to include FLOPS:
    minimize total compute subject to latency and
    memory budgets. Most common for
    throughput-sensitive workloads.
\end{itemize}

The two reduce to each other by adding the other
resource to the budget vector, so without loss of
generality we take $\phi(\bar{c}) = \bar{c}_{\LAT}$
below.

\subsection{The dual}

The Lagrangian for the primal with $\phi = \bar{c}_{\LAT}$
and budget constraints $\bar{c}_r \leq B_r$ for
$r \in \{\VRAM, \SMEM, \REG\}$ (we omit $\LAT$ from
the constraint side since it is the objective) is
\begin{equation}\label{eq:lagrangian}
  \mathcal{L}(P, \lambda_V, \lambda_S, \lambda_R) =
  \bar{c}_{\LAT}(P) +
  \lambda_V (\bar{c}_{\VRAM}(P) - B_{\VRAM}) +
  \lambda_S (\bar{c}_{\SMEM}(P) - B_{\SMEM}) +
  \lambda_R (\bar{c}_{\REG}(P) - B_{\REG}).
\end{equation}

The dual function is
\begin{equation}
  g(\lambda_V, \lambda_S, \lambda_R) =
  \min_{P \in \mathrm{plans}_G(W)}
  \mathcal{L}(P, \lambda_V, \lambda_S, \lambda_R),
\end{equation}
and the dual problem is $\max_{\lambda \geq 0} g(\lambda)$.

For fixed $\lambda$, $\mathcal{L}$ is a weighted sum over
$\bar{c}$, and the inner minimization reduces to a
single unweighted chart parse with the combined scalar
cost $\phi_\lambda(\bar{c}) = \bar{c}_{\LAT} +
\sum_r \lambda_r \bar{c}_r$. We call this the
\emph{weighted parse at Lagrangian price $\lambda$}.

\subsection{Subgradient update}

The outer maximization over $\lambda$ is solved by
subgradient ascent:

\begin{definition}[Subgradient update]\label{def:subgrad}
  Starting from $\lambda^{(0)} = 0$:
  \begin{enumerate}[nosep,label=(\arabic*)]
    \item Solve the weighted parse at $\lambda^{(t)}$,
      obtaining $P^{(t)}$.
    \item Compute the slack vector $s^{(t)} =
      \bar{c}(P^{(t)}) - B$ (restricted to the
      constraint dimensions).
    \item Update: $\lambda^{(t+1)} = \max(0,
      \lambda^{(t)} + \eta_t \cdot s^{(t)})$
      componentwise, with step size $\eta_t
      = c / \sqrt{t}$.
  \end{enumerate}
\end{definition}

Convergence of this subgradient scheme under the LP
relaxation is standard (Bertsekas 1999, Chapter 6).

\subsection{Dual prices as diagnostics}

\begin{proposition}[Dual price interpretation]
  \label{prop:dual-prices}
  At the optimal dual solution $\lambda^*$:
  \begin{itemize}[nosep]
    \item $\lambda_r^* > 0$ iff the constraint on
      resource $r$ is active (binding) at the primal
      optimum;
    \item $\lambda_r^*$ is the marginal decrease in
      the objective per unit increase in the budget
      $B_r$.
  \end{itemize}
\end{proposition}

\begin{proof}
  Complementary slackness of LP duality: at optimum,
  either $\lambda_r^* = 0$ or $s_r^* = 0$, and the
  sensitivity interpretation is the standard envelope
  theorem applied to $\mathcal{L}$.
\end{proof}

In operational terms: if the planner returns a plan
whose dual price $\lambda_{\REG}^*$ is high, the user
knows that loosening the register budget (e.g.\ by
using a device with a larger register file) would
most improve execution latency. This is a direct,
actionable architectural insight.

\subsection{Achievability and the integrality gap}

\begin{theorem}[Lagrangian achievability]
  \label{thm:lagrangian}
  \footnote{Class \textsf{AH}, status \omark\ (open
  integrality gap). Warrant: LP duality applies to
  the LP relaxation of the plan-selection problem.
  Qualifier: the gap between LP-optimal and
  IP-optimal plan cost is bounded only empirically
  in our experiments; a formal worst-case bound is
  open (\S1 item X3). Finite-$\kappa$ valid.}
  The LP relaxation of the primal problem has a
  computable dual optimum $\lambda^*$, and the
  primal plan obtained by executing a chart parse at
  price $\lambda^*$ is within a bounded factor of
  the LP-relaxed optimum. The gap is zero when the
  grammar's production costs are affine in batch
  size.
\end{theorem}

\begin{proof}
  See Appendix \ref{app:b}, Subsection
  \ref{app:b-lagrangian}. The affine-in-batch case
  is classical; the general case depends on the
  step function character of certain productions
  (e.g.\ batch saturation of LCM packing) and is
  bounded by a small constant (typically $\leq 2$)
  in practice.
\end{proof}

\subsection{Practical consequences}

Three consequences worth stating explicitly:

\begin{itemize}[nosep]
  \item The planner can report, alongside each plan,
    the three dual prices $(\lambda_{\VRAM}^*,
    \lambda_{\SMEM}^*, \lambda_{\REG}^*)$. These form
    a human-readable diagnostic of which resource is
    the current bottleneck.
  \item Tightening any single budget produces a new
    plan whose cost may differ by at most the
    corresponding $\lambda_r^*$ per unit of tightening.
    This is exactly the sensitivity the user wants
    when deciding which resource to invest in.
  \item The weighted parse at a fixed $\lambda$ is a
    single Earley pass; the outer subgradient
    iteration typically converges in $O(\log 1 /
    \epsilon)$ iterations for tolerance $\epsilon$.
    Total cost is therefore the Earley complexity of
    Theorem \ref{thm:complexity} multiplied by a
    logarithmic factor.
\end{itemize}

\subsection{Relationship to constrained parse-forest optimization}

The Lagrangian dual formulation is a specific instance
of the general framework of constrained CFG parsing
(Chiang et al.\ 2013; Lopez 2008; Huang 2008). The
specialization to our setting is that the resource
vector and its associated cost lattice are
multi-dimensional, and the productions combine costs
by different rules per dimension (sum, max). These
additions do not break the LP-duality argument; they
only complicate the Lagrangian gradient, which remains
a subgradient at the piecewise-linear points where
two productions' costs are equal.
